\documentclass[11pt]{article}
\usepackage{fontspec,amsmath,amssymb,amsthm,hyperref}
\usepackage[margin=1in]{geometry}
\setlength{\emergencystretch}{3em}
\title{An abstract family of compatible finite-coupling references}
\author{LangevinLean mathematical exposition}
\date{}
\begin{document}\maketitle
\noindent Source module: \texttt{UnadjustedLangevin/ReferenceFamily.lean}. Compile this proof document with XeLaTeX or LuaLaTeX. The accompanying HTML pairs the mathematical argument with the exact Lean statements and records the scope of the formalization.


This interface isolates the passage from a bound for each initial coupling to a theorem with the exact initial Wasserstein distance. A concrete product construction later supplies this interface without assuming an optimal plan.


\section{FiniteCoupling}


Restrict the coupling type to those with \(\operatorname{couplingCost}\ne\infty\). This subtype records the integrability needed to realize initial class distances and convert them to real values while taking the Wasserstein infimum.



\noindent Lean declarations: \texttt{FiniteCoupling}.

\section{ReferenceFamily}


For each finite coupling, supply initial and noise \(L^2\) variables and an \texttt{L2Reference}. The fields identify initial law, bound initial class distance by that coupling's cost, and require one common output law across the family. These are explicit construction obligations, with no convergence estimate among them.



\noindent Lean declarations: \texttt{ReferenceFamily}.

\section{ReferenceFamily.output finite}


Choose one finite initial coupling using the finite initial Wasserstein premise. Its supplied numerical/reference \(L^2\) variables yield a finite-cost output coupling. The output-law and stationary-reference fields identify its marginals, proving \(W_2(\nu_n,\pi)<\infty\).



\noindent Lean declarations: \texttt{ReferenceFamily.output\_finite}.

\section{Pass the analytic bound through the infimum}


Apply the analytic \(L^2\) estimate to every finite coupling realization in the family. Replace its initial error by the coupling cost and its numerical law by the common \(\nu_n\). Because \((1-mh)^n\ge0\), the affine finite-coupling infimum theorem gives the displayed bound with the true initial Wasserstein distance.



\noindent Lean declarations: \texttt{dalalyan\_small\_step\_of\_reference\_family}.

\end{document}
